% chapters/chapter1/sec1.tex

\section{The Essence, Evolution, and Core Characteristics of Integrated Sensing and Communications}
\label{sec:ch1-essential-definition}

% TODO: This section is drafted by Huixin Man.
% Writing objectives:
% 1. Introduce the developmental background of wireless communication systems;
% 2. Motivate the 6G demand for the convergence of communications, sensing, and intelligence;
% 3. Explain the research significance of ISAC.

\subsection{Why Is It Necessary to Redefine Integrated Sensing and Communications}
\label{subsec:ch1-why-redefine-isac}
\subsubsection{The Traditional Separation Between Communication and Radar Systems}
\label{subsubsec:ch1-traditional-separation}

Wireless communication systems and radar systems are two of the most important branches of modern electronic information technology. From the perspective of engineering history, although communication systems and radar systems both belong to electromagnetic information systems and both rely on the transmission, propagation, and reception of wireless signals, they have served fundamentally different tasks since their inception. As a result, they have evolved along two relatively independent technological paths. Since the emergence of radio communication at the end of the nineteenth century, the core objective of wireless communication systems has remained unchanged: to realize reliable and efficient information transmission between a transmitter and a receiver. From the viewpoint of classical information theory, under the Additive White Gaussian Noise (AWGN) channel model, Shannon's theory reveals the fundamental relationship among channel capacity, bandwidth, and Signal-to-Noise Ratio (SNR) \cite{Shannon1948}:
\begin{equation}
C = B \log_2 \left( 1 + \frac{S}{N} \right)
\end{equation}
where $C$ is measured in bit/s, $B$ denotes the bandwidth (Hz), and $S$ and $N$ represent the average powers (W) of the signal and noise, respectively. This formula highlights two central goals of communication system design: maximizing the information transmission rate and ensuring transmission reliability under a given SNR condition.

Driven by these goals, communication systems have undergone more than a century of continuous evolution, from analog to digital, from narrowband to broadband, from single-antenna to multi-antenna, and from single-user to massive multiuser systems \cite{Goldsmith2005,Andrews2014}. From the analog cellular networks of the First-Generation Mobile Communication System (1G) to the three representative service scenarios of the Fifth-Generation Mobile Communication System (5G), namely Enhanced Mobile Broadband (eMBB), Ultra-Reliable Low-Latency Communications (URLLC), and Massive Machine-Type Communications (mMTC), the design methodology of communication systems has consistently revolved around source coding, channel coding, modulation and demodulation, multiple access, and resource scheduling \cite{TS38300}. In communication systems, the wireless channel is viewed as a ``pipe'' for information delivery, while its time-varying characteristics, such as multipath fading and Doppler spread, are usually regarded as adverse effects that must be mitigated or compensated for \cite{Goldsmith2005}.

Almost in parallel with wireless communication technology, Radio Detection and Ranging (RADAR) technology began to emerge in the early twentieth century and developed rapidly during World War II \cite{Skolnik2008}. In sharp contrast to communication systems, the core objective of radar/sensing systems is not information transmission, but the extraction of physical information about targets or the surrounding environment from echo signals generated by electromagnetic interactions with objects and scenes \cite{Skolnik2008,Richards2014}. Specifically, the typical tasks of radar systems include:
\begin{itemize}
    \item \textbf{Target detection:}  
    determining whether a target of interest exists within a specified spatial region, which is a classical binary hypothesis testing problem;
    \item \textbf{Parameter estimation:}  
    estimating target parameters such as range, velocity, and angle, whose theoretical lower performance bounds are often characterized by the Cram\'er--Rao Lower Bound (CRLB) \cite{Kay1993};
    \item \textbf{Target recognition and classification:}  
    identifying target categories by analyzing features of the echo signal, such as micro-Doppler signatures;
    \item \textbf{Environmental imaging and reconstruction:}  
    for example, Synthetic Aperture Radar (SAR) and Inverse Synthetic Aperture Radar (ISAR) achieve fine-grained characterization of targets or terrain through high-resolution imaging \cite{Richards2014}.
\end{itemize}

In radar systems, transmit signal design aims at achieving favorable ambiguity function characteristics, namely, a sharp mainlobe and low sidelobes in the joint range--Doppler plane, so as to enable high-resolution estimation of target parameters \cite{Levanon2004,Richards2014}. This design philosophy is fundamentally different from that of communication systems, which emphasizes maximizing spectral efficiency and minimizing bit error rate \cite{Liu2020JointRadarComm,Zhang2021ISAC6G}.

Because communication systems and radar systems differ fundamentally in their core objectives, they have historically followed very different development paths with respect to spectrum usage, hardware platforms, signal design, and receiver processing \cite{Liu2020JointRadarComm,Liu2022ISAC}.
\begin{enumerate}[label=(\arabic*)]
    \item \textbf{Spectrum usage}
    
    Communication systems and radar systems have long operated over separately allocated spectrum resources. Under the historical spectrum management framework, the International Telecommunication Union (ITU), through the \emph{Radio Regulations} and the World Radiocommunication Conference (WRC), allocates dedicated frequency bands to different services. For example, cellular mobile communication systems mainly operate in bands such as 700 MHz to 3.5 GHz, while 5G New Radio (5G NR) further extends into the millimeter-wave range. In contrast, radar systems are distributed over bands such as the S band (2--4 GHz), C band (4--8 GHz), and X band (8--12 GHz). Such physical separation of the spectrum has enabled the two types of systems to operate independently with limited mutual interference.
    
    \item \textbf{Hardware platforms}
    
    Communication devices, such as base stations and user terminals, differ significantly from radar equipment in hardware architecture. The Radio Frequency (RF) front end of communication systems is typically optimized for the transmission and reception of linearly modulated signals, with emphasis on linearity and efficiency. By contrast, radar systems, especially pulse radar systems, often employ high-power amplifiers at the transmitter, while the receiver uses low-noise amplifiers and Analog-to-Digital Converters (ADCs) with high dynamic range to deal with weak echoes over a large dynamic range. Antenna design also differs substantially. Military radar systems often employ high-gain directional antennas or phased arrays with narrow beams, whereas cellular base stations usually require omnidirectional or sectorized antennas to balance coverage requirements.
    
    \item \textbf{Signal design}
    \begin{sloppypar}
    Communication signal design centers on carrying modulated information. Typical communication signals, such as Orthogonal Frequency-Division Multiplexing (OFDM), achieve efficient information transmission by loading independent data symbols onto frequency-domain subcarriers. Their waveform design focuses on key metrics such as spectral efficiency, Peak-to-Average Power Ratio (PAPR), and out-of-band leakage \cite{Goldsmith2005,Andrews2014}. In contrast, radar signal design aims to optimize sensing performance. Classical radar signals include Linear Frequency-Modulated (LFM) pulses and phase-coded pulses, whose design criteria revolve around the shape of the ambiguity function to simultaneously achieve high range and velocity resolution with low sidelobes \cite{Levanon2004}. Radar signals are typically deterministic and known to the receiver, since matched filtering requires prior knowledge of the transmitted waveform. By contrast, from the radar perspective, communication signals are usually ``random'' because the data symbols are not known a priori to the sensing receiver \cite{Liu2020JointRadarComm,Zhang2021ISAC6G}.
    \end{sloppypar}

    \item \textbf{Receiver processing}
    
    The core task of a communication receiver is signal detection and decoding, namely, recovering the information bits transmitted by the sender in the presence of noise and interference. A typical processing chain includes synchronization, channel estimation, equalization, demodulation, and decoding \cite{Goldsmith2005}.
    
    By contrast, the core task of a radar receiver is target detection and parameter extraction. Its typical processing chain includes matched filtering (pulse compression), Constant False Alarm Rate (CFAR) detection, Doppler processing, and angle estimation. The focus of a radar receiver is to extract physical target parameters from the echo signal rather than to recover modulated information \cite{Skolnik2008,Richards2014}.
\end{enumerate}

\begin{table}[!ht]
\centering
\renewcommand{\arraystretch}{1.2}
\begin{tabularx}{\textwidth}{|>{\raggedright\arraybackslash}p{3.2cm}|
>{\raggedright\arraybackslash}X|
>{\raggedright\arraybackslash}X|}
\hline
\textbf{Comparison Dimension} & \textbf{Communication Systems} & \textbf{Radar/Sensing Systems} \\
\hline
\textbf{Core objective} & Reliable information transmission & Target detection and parameter estimation \\
\hline
\textbf{Spectrum usage} & Licensed communication bands & Dedicated radar bands \\
\hline
\textbf{Signal characteristics} & Random data-bearing modulation & Deterministic known waveform \\
\hline
\textbf{Role of the channel} & A ``pipe'' for information delivery, to be overcome & A ``carrier'' of sensing information, to be analyzed \\
\hline
\textbf{Transmit signal design criterion} & Spectral efficiency, bit error rate & Ambiguity function, resolution \\
\hline
\textbf{Receiver processing objective} & Bit recovery & Parameter extraction \\
\hline
\textbf{Performance metrics} & Capacity, throughput, latency & Detection probability, estimation accuracy \\
\hline
\textbf{Typical hardware architecture} & Base station/terminal & Radar station/sensor \\
\hline
\end{tabularx}
\caption{Core comparison between communication systems and radar/sensing systems}
\label{tab:comm-vs-radar}
\end{table}

Table~\ref{tab:comm-vs-radar} summarizes the major differences between traditional communication systems and radar/sensing systems across multiple dimensions. It should be noted that their independent evolution was historically reasonable through the 5G era and earlier generations. First, from the spectrum perspective, communication and radar systems in the sub-6 GHz range were largely physically separated, and spectrum conflict was not yet prominent. Second, from the service perspective, mobile communication networks before 5G mainly served human-to-human communication such as voice and data, and did not require explicit sensing or understanding of the physical environment. Third, from the perspective of technological maturity, the communication and radar technology stacks had both become highly specialized, and further deepening each stack independently was an efficient path.

However, this separated paradigm is now revealing increasingly severe limitations. First, spectrum scarcity is aggravating the contradiction. As mobile communications expand into the millimeter-wave (mmWave) and terahertz (THz) bands, the operating frequency ranges of communication and radar systems are becoming adjacent or even overlapping, placing increasing pressure on the traditional static allocation paradigm \cite{TS38104,Liu2022ISAC}. For example, the 24.25--52.6 GHz mmWave range of 5G NR is spectrally close to multiple radar applications, such as automotive radar and weather radar \cite{Zheng2019RadarCommunicationCoexistence}. The scarcity of spectrum makes it increasingly unsustainable for communication and radar systems to occupy their own dedicated bands for long periods.

Second, emerging applications require joint communication-and-sensing capability. Autonomous vehicles must communicate with roadside infrastructure while simultaneously perceiving their surroundings; unmanned aerial vehicles (UAVs) must maintain control links while performing obstacle sensing; and large numbers of Internet of Things (IoT) devices in smart cities may need to upload data while also sensing their environment \cite{Liu2022ISAC,De2021convergent}. In such scenarios, separately deploying communication and sensing functionalities often leads to higher system complexity, cost, and latency.

Finally, independent deployment causes increasingly significant resource waste. Communication base stations and radar systems both possess RF front ends, antenna arrays, and signal processing units, and these resources have considerable sharing potential at both the physical and platform levels \cite{Liu2020JointRadarComm,Liu2022ISAC}. Independent deployment therefore results in duplicated investments in hardware, energy consumption, and site infrastructure. These limitations jointly constitute the fundamental motivation for the development of Integrated Sensing and Communications (ISAC) technology and provide the historical background for the subsequent discussion of communication--sensing convergence requirements.

\subsubsection{The Fundamental Demand for the ``Convergence of Communication and Sensing'' in New Wireless Network Scenarios}
\label{subsubsec:ch1-fundamental demands}

Fundamentally, the emergence of ISAC is not merely the result of a new technical concept, but a response to a structural shift in the capability requirements of next-generation wireless network applications. Future wireless networks are no longer expected only to provide traditional ``connectivity''; they must also possess ``environmental sensing capability'' and ``task support capability.'' In this context, communication systems are evolving from pure information transmission platforms into comprehensive information infrastructures that integrate sensing, cognition, and collaborative decision-making \cite{Liu2022ISAC}.

First, spectrum scarcity and the duplicated occupancy of resources by separate communication and sensing systems are important practical drivers behind communication--sensing integration. Spectrum is the fundamental physical resource shared by both wireless communication systems and radar/sensing systems. As global mobile data traffic continues to grow, the demand for spectrum keeps increasing, and the scarcity of high-quality spectrum, especially contiguous low- and mid-frequency bands, is becoming more pronounced. According to the frequency allocation and regulatory documents released by the ITU and national spectrum regulatory authorities, within the sub-6 GHz range, mobile communications, satellite services, radiolocation, aeronautical and maritime radar, and other dedicated wireless services already occupy a large amount of spectrum, leaving fewer and fewer high-quality contiguous bands for exclusive use by new systems \cite{ITU_R_RR2024,miit_radio_freq_2026_en}. Meanwhile, mobile communication systems are extending further into the mmWave and even THz bands, which inevitably leads them into frequency ranges historically used by radar, remote sensing, and radiolocation systems, or into bands adjacent to them. Around the 24 GHz band, for example, 5G mmWave services have expanded above 24.25 GHz, while the adjacent 23.6--24.0 GHz band has long been used for passive Earth exploration satellite sensing; moreover, the 24 GHz band has historically been widely used for short-range automotive radar. This indicates that as communication systems continue to move toward higher frequencies, the pressure for adjacent-channel coexistence or even co-channel coordination with traditional sensing and remote-sensing systems will increase significantly. Against this background, continuing to allocate separate spectrum resources to communication and sensing systems is becoming increasingly difficult and economically inefficient in terms of overall spectrum utilization. One of the core advantages of ISAC is that it can fundamentally alleviate the contradiction between spectrum scarcity and duplicated system occupancy by enabling communication and sensing to share spectrum, infrastructure, and signal processing chains. More importantly, if a unified waveform can be designed to simultaneously carry communication information and support environmental sensing, the resulting gain can exceed that of simple time-sharing or frequency-sharing, thereby bringing additional spectral efficiency improvements \cite{Liu2022ISAC,Hassanien2015dual}.

Second, the demand for \emph{native sensing capability} at terminals, base stations, and edge nodes is increasing. In traditional network architectures, sensing is usually handled by independent sensor systems, such as radar, cameras, or Light Detection and Ranging (LiDAR), while communication networks are mainly responsible for data transfer; the coupling between the two is therefore relatively weak. However, many 6G-oriented scenarios are pushing network nodes themselves to possess built-in sensing capability, meaning that sensing is no longer an external peripheral module attached to the communication system, but is gradually becoming a basic function of the network node itself \cite{Zhang2021ISAC6G}. This trend is particularly evident in several representative scenarios. For example, in vehicular networks and autonomous driving, in-vehicle communication modules already possess strong RF, antenna, and baseband processing capabilities. If the communication link can additionally support target detection, localization, and environmental sensing, dependence on separate sensing hardware may be reduced, thereby lowering system complexity and deployment cost. In low-altitude economy and UAV networks, aircraft often need to simultaneously complete communication, localization, obstacle avoidance, and environmental sensing under tight payload and power constraints; embedding sensing capability into the communication link can improve integration and energy efficiency. In smart factories and industrial IoT scenarios, massive numbers of devices require not only reliable connectivity but also real-time awareness of personnel, equipment, and material states; if the communication infrastructure can simultaneously provide sensing services, the large-scale deployment of additional sensors can be significantly reduced. In smart city and public safety applications, densely deployed base stations that can also perform traffic monitoring, anomaly detection, and spatial state sensing would substantially expand the functional boundaries and service value of the existing network infrastructure.

At a deeper level, the necessity of ISAC stems from a fundamental redefinition of the role of future wireless networks. In 5G and earlier generations, the core task of wireless networks was mainly to provide efficient and reliable information transmission; the network was essentially an ``information transport pipe,'' and its key concerns were rate, latency, coverage, and connection density, with relatively limited involvement in understanding the physical environment itself. However, in the 6G vision, wireless networks are gradually transforming from traditional ``connectivity platforms'' into integrated service platforms that combine connectivity, sensing, computation, and control \cite{Liu2022ISAC,Zhang2021ISAC6G}. This transformation implies that the network must not only transmit information, but also understand the physical environment in which it operates, such as sensing user location, velocity, and activity state, characterizing the distribution of surrounding objects, and obtaining spatial structural information about the channel. The network must not only provide links, but also support intelligent decision-making, for example by jointly optimizing beam management, resource scheduling, mobility management, and task allocation using both communication state information and environmental sensing results. For highly reliable and low-latency closed-loop applications such as autonomous driving, industrial control, and remote operation, the network must further support end-to-end coordination from environmental sensing to information transmission, edge computing, and control execution, where sensing information is often the starting point of the control loop. For this reason, sensing is no longer an optional auxiliary feature of future wireless networks, but is becoming an essential component of their native capability set.

Once these requirements are recognized, a natural question arises: if future networks need sensing capability, why not simply adopt an ``external'' architecture consisting of a communication system plus an independent sensing system, rather than emphasizing integrated design? In fact, such an external approach is already common in practice, for example in road infrastructure where communication base stations and independent radar units are co-deployed, or in intelligent vehicles equipped with both communication modules and mmWave radar. However, as application scenarios impose increasingly stringent requirements on real-time operation, collaboration, system cost, and deployment scale, the limitations of this separated architecture become more evident. First, independent communication and sensing systems often require additional data interfaces, synchronization mechanisms, and information exchange processes, which introduce non-negligible processing overhead and transmission delay. This is particularly undesirable in latency-sensitive scenarios such as autonomous driving and industrial control. Second, when communication and sensing systems are designed and optimized separately, they lack deep coordination in waveform structures, resource allocation, air-interface protocols, and hardware platforms, making it difficult to achieve globally optimal resource utilization and performance balance. Third, deploying both communication equipment and independent sensing equipment at every node incurs higher hardware cost, power consumption, site cost, and maintenance cost, which is economically unattractive for large-scale commercial deployment. Fourth, communication data and sensing data generated by different systems and standards often differ in timestamping, coordinate systems, data formats, and processing pipelines, making subsequent fusion more complex and system implementation more challenging \cite{Liu2022ISAC,Hassanien2015dual}. By contrast, ISAC improves spectrum utilization fundamentally, enhances communication--sensing synergy, and provides a unified technical pathway toward higher integration, lower cost, and stronger task support capability through integrated design of architecture, signal processing, resource scheduling, and receiver processing.

\subsubsection{ISAC Cannot Be Simply Understood as a Combination of ``Communications + Radar''}
\label{subsubsec:ch1-isac-not-simple-combination}

After clarifying the necessity of ISAC, it is also important to dispel a common conceptual misunderstanding: ISAC is \emph{not} simply a matter of attaching a communication system and a radar system together. In existing literature and engineering practice, there are multiple levels at which communications and sensing can be combined, ranging from mere physical co-location, to hardware sharing, to spectrum coexistence, and then to deep integration at the waveform, protocol, and network architecture levels. These schemes differ fundamentally in their \emph{degree of integration} \cite{Liu2022ISAC,Zhang2021ISAC6G}. Without carefully distinguishing among these notions, one may mistakenly regard many systems in which communication and sensing merely coexist as ISAC systems, thereby obscuring the true technical meaning of ISAC.

The most naive interpretation is that a communication device and a radar device are installed at the same physical site, each independently performing its own function without any substantial system-level coordination. This realizes only \emph{spatial co-location}. Its main advantages are deployment convenience and partial reuse of infrastructure such as sites and power supply. However, from the system perspective, it still consists of two independent devices operating in parallel. Communication and sensing remain separate in signal design, resource scheduling, information processing, and performance objectives, and therefore this arrangement cannot be regarded as a genuinely integrated system.

A further step beyond physical co-location is to reuse part of the hardware platform, for example by sharing the RF front end, antenna array, or baseband processing platform, while switching between communication and sensing functions over different time slots, frequency bands, or operating modes. Compared with completely independent deployment, such solutions can reduce device size, manufacturing cost, and energy consumption, and they do reflect an initial level of coordination between communication and sensing. Nevertheless, communication and sensing tasks are usually still orthogonally separated in time or frequency, rather than jointly operating under a unified signal, protocol, or resource optimization framework. Hence, such schemes are better viewed as \emph{shared-platform implementations} of communications and sensing, or as primitive forms of ISAC, rather than a complete realization of the ISAC concept \cite{Hassanien2015dual}.

Another class of systems that is easily confused with ISAC is \emph{spectrum coexistence}. In this class of systems, communication and radar/sensing systems share the same spectrum resources and employ techniques such as cognitive access, dynamic spectrum allocation, interference suppression, beam isolation, or interference alignment to reduce their mutual impact \cite{Liu2020JointRadarComm}. The focus of spectrum coexistence is the following question: when two originally independent systems are forced or allowed to operate in the same spectrum, how can they be made to interfere with each other as little as possible? Therefore, the core issue in coexistence is compatibility and interference management, rather than unified design and synergistic enhancement of functionality. In other words, under the coexistence framework, communication and sensing remain two independent systems with separate objectives, architectures, and processing chains, even though they share the same physical spectrum. This remains fundamentally different from the ISAC philosophy of simultaneously realizing communication and sensing \emph{within one unified system} \cite{Liu2022ISAC}.

In a stricter sense, true ISAC aims at \emph{deep integration} of communication and sensing across multiple layers of the system. Such integration is not limited to equipment sharing, but manifests itself at the levels of signals, resources, hardware, protocols, networks, and service logic. Its objective is no longer merely to avoid conflict, but to achieve synergistic gains in communication and sensing performance through joint design \cite{Liu2022ISAC,Zhang2021ISAC6G}. Specifically, this deep integration is mainly reflected in the following aspects:

\begin{itemize}
    \item \textbf{Signal-level integration:} the same waveform simultaneously carries communication information and enables environmental sensing, rather than simply superimposing communication and sensing signals or alternating between them. Ideally, the communication data itself becomes part of the sensing waveform design, while the environmental information contained in sensing echoes can in turn support communication tasks such as link estimation, beam management, or channel prediction.

    \item \textbf{Resource-level integration:} wireless resources in the time, frequency, space, and power domains are jointly occupied by communication and sensing tasks and allocated under a unified optimization framework. Resource management is therefore no longer performed separately for ``the communication system'' and ``the sensing system,'' but must jointly consider tradeoffs among communication rate, latency, and reliability on the one hand, and sensing detection probability, estimation accuracy, and resolution on the other.
    
    \begin{sloppypar}
    \item \textbf{Hardware-level integration:} antenna arrays, RF chains, analog-to-digital/digital-to-analog converters, and baseband processing units are genuinely shared between communication and sensing functions, and the hardware architecture is designed from the outset to satisfy the requirements of both tasks, rather than merely running two independent software stacks on the same platform in alternation.
    \end{sloppypar}

    \item \textbf{Protocol-level integration:} air-interface protocols, frame structures, pilot and reference signal design, feedback mechanisms, random access procedures, and link maintenance processes must natively support sensing, rather than simply adding a sensing module on top of a conventional communication protocol. In other words, sensing should be incorporated into the protocol design from the beginning.

    \item \textbf{Network-level integration:} sensing must be incorporated into the overall network architecture, including sensing task triggering and scheduling, sensing data aggregation and distribution, multi-node cooperative sensing, and feedback from sensing results to network resource management. In this sense, sensing is no longer a local function of an individual node but part of a network-wide capability.

    \item \textbf{Service-level integration:} communication services and sensing services jointly support upper-layer applications, forming a complete service loop from environmental sensing to information transmission, then to computing, decision-making, and control execution. In scenarios such as autonomous driving, intelligent manufacturing, low-altitude coordination, and smart cities, this service-layer closed loop is one of the key manifestations of ISAC value.
\end{itemize}

\begin{sloppypar}
Therefore, the most fundamental distinction between ISAC and a traditional combined ``communication system + radar system'' solution does not lie in whether the two are installed at the same location or whether some hardware resources are shared. Rather, it lies in whether communication and sensing are jointly designed and jointly optimized across layers under a unified system objective. Only when communication and sensing are no longer treated as separate auxiliary functions, but instead become organically unified across waveforms, resources, protocols, networks, and service workflows, does a system truly exhibit the essential characteristics of ISAC. This multilayer, deeply coupled, task-oriented integration mechanism forms the central theme of the remainder of this book.
\end{sloppypar}

\subsection{The Basic Definition and Conceptual Boundaries of ISAC}
\label{subsec:ch1-basic-definition}
\subsubsection{The Definition of ISAC}
\label{subsubsec:ch1-definition-isac}

Having discussed the practical motivations behind ISAC, its core idea, and its distinction from concepts such as simple co-location or spectrum coexistence, we now proceed to provide a basic definition of ISAC and clarify its conceptual boundaries. It should be emphasized that although academia and industry still differ somewhat in their preferred implementation paths for ISAC, there is broad consensus on its essential meaning: the core of ISAC lies not in mechanically superimposing communication and sensing functions, but in achieving their joint design, coordinated operation, and bidirectional empowerment under a unified system objective \cite{Liu2022ISAC,Liu2020JointRadarComm}.

In this textbook, ISAC is defined as follows:
\begin{quote}
\textbf{ISAC refers to a class of technologies in which, under a unified system framework, communication and sensing functionalities are simultaneously enabled through shared or jointly designed signals, hardware, spectrum, air-interface mechanisms, and network mechanisms.}
\end{quote}

This definition contains several key elements, each of which deserves further explanation.
\begin{enumerate}[label=(\arabic*)]
    \item \textbf{A unified system framework}
    
    The primary feature of ISAC is \emph{unity}: communication and sensing are not loosely stitched together as two independent subsystems, but are modeled, designed, and operated within a common system framework \cite{Liu2022ISAC,Zhang2021ISAC6G}. This unified framework is often reflected in multiple aspects, including a unified signal processing chain, unified wireless resource management, unified network architecture, and a unified performance evaluation perspective. It should be stressed that such ``unity'' does not imply that communication and sensing have identical physical goals, performance metrics, or processing methods. On the contrary, their requirements for rate, latency, reliability, detection probability, estimation accuracy, and resolution are often different or even conflicting. The meaning of ``unity'' is that these differences are not treated separately and optimized in isolation, but are jointly considered and coordinated within a common framework from the outset.
    
    \item \textbf{Sharing or co-design}
    
    In ISAC systems, the integration between communication and sensing is typically manifested in two basic forms: \emph{sharing} and \emph{co-design}, with practical systems usually lying on a continuum between them \cite{Liu2020JointRadarComm,Gonzalez2024integrated}. In the sharing mode, communication and sensing may use the same waveform, spectrum resources, hardware platform, or processing modules. For example, the same OFDM waveform may both transmit information to users and provide environmental sensing information through analysis of propagation paths or echo structures. In the co-design mode, communication and sensing do not necessarily use exactly the same resources or signal forms, but they are jointly considered and jointly optimized at the system design stage. For example, beamforming or power allocation may need to simultaneously balance communication coverage and data rate against sensing directionality and estimation precision. Thus, ISAC is not restricted to the extreme case where all resources are fully shared; more generally, it emphasizes the intrinsic coordination between communication and sensing in design objectives and resource utilization.
    
    \item \textbf{Simultaneous communication and sensing functionalities}
    
    The term ``simultaneous'' in the ISAC definition should be understood from both the system objective and operational mechanism perspectives. First, at the system objective level, communication and sensing are treated as two equally important core functions from the design stage, rather than as a primary function with the other serving merely as an add-on \cite{Liu2022ISAC}. Second, at the operational level, ``simultaneous'' may mean strict parallel operation, where the same signal or resource block performs both communication and sensing at the same instant, or it may refer to very fast switching and quasi-simultaneous execution on a short time scale. In the latter case, even if communication and sensing are not exactly coincident at every instant, the system still falls within the ISAC category as long as the two functions are modeled on an equal footing and jointly managed in terms of system logic, resource control, and task scheduling. In other words, whether a system belongs to ISAC does not hinge on absolute temporal coincidence, but on whether communication and sensing are jointly designed and jointly operated within a unified framework.
    
    \item \textbf{The broad scope of sensing objects}
    
    In the ISAC context, ``sensing'' should be interpreted in a broad sense, far beyond the point targets or airborne targets commonly associated with traditional radar \cite{Liu2022ISAC,Gonzalez2024integrated}. More generally, the sensing objects in ISAC may include at least the following categories: first, the external physical environment, such as buildings, walls, road structures, terrain, and even weather conditions; second, moving targets, such as vehicles, pedestrians, and UAVs, including their detection, localization, tracking, and identification; third, network objects, such as user location, mobility state, spatial distribution, and behavioral patterns; fourth, devices and terminals, including device presence detection, state estimation, and activity monitoring; fifth, the electromagnetic environment, such as propagation paths, scatterer distributions, angle--delay structures, and interference states; and sixth, vital signs and human activities, such as respiration monitoring, heartbeat estimation, fall detection, gesture recognition, and human activity recognition. It is precisely because sensing in ISAC has such a broad meaning that the application boundary of ISAC extends far beyond that of traditional radar and becomes closely connected to future wireless networks, ambient intelligence, and human--machine interaction.
    
    \item \textbf{Bidirectional empowerment between communication and sensing}
    
    Another essential feature of ISAC is that the relationship between communication and sensing is not one-directional, but rather one of \emph{bidirectional empowerment} \cite{Liu2022ISAC,Zhang2021ISAC6G}. On the one hand, sensing can assist communication: by acquiring user position, target motion states, obstacle distributions, spatial scattering structures, or link visibility information, sensing can provide prior information for beam management, link establishment, channel prediction, handover decisions, and resource scheduling, thereby improving communication reliability, stability, and resource efficiency. On the other hand, communication can empower sensing: communication systems naturally provide RF signal sources, network coverage, time synchronization, distributed node deployment, and feedback mechanisms, all of which serve as natural carriers and infrastructure support for environmental sensing, target detection, and multi-node cooperative sensing. Especially in large-scale deployments such as cellular networks, Wi-Fi networks, and vehicular networks, the distributed and continuously online nature of communication infrastructure makes it more suitable than traditional standalone radar for wide-area sensing platforms. It is this dual mechanism of ``sensing-assisted communication'' and ``communication-empowered sensing'' that makes ISAC fundamentally different from simply adding a sensing module to a communication system. Instead, it should be viewed as a unified technological framework that achieves a leap in overall system performance through the fusion of two functionalities.
\end{enumerate}

Taken together, the basic definition of ISAC contains at least three core implications: first, communication and sensing are jointly designed and jointly operated at the system level; second, communication and sensing not only share resources, but can also create synergistic gains through joint optimization; and third, sensing in ISAC is broad, networked, and capable of feeding back to optimize communication performance. Understanding these key elements will help us more accurately grasp the research scope and technical boundaries of ISAC in subsequent discussions of signal models, system architectures, and performance tradeoffs \cite{Liu2022ISAC,Liu2020JointRadarComm,Gonzalez2024integrated}.

\subsubsection{Boundary Distinctions Between ISAC and Related Concepts}
\label{subsubsec:ch1-boundary-distinction}

In academic research, engineering practice, and standardization discussions surrounding ISAC, several closely related but non-equivalent technical concepts frequently appear. Accurately distinguishing the connections and differences between ISAC and these related concepts is important for building a clear technical picture, avoiding terminological confusion, and correctly understanding the actual scope of different research efforts. Broadly speaking, these related concepts may be grouped into three categories: first, the \emph{coexistence} problem, which focuses on interference management; second, \emph{joint radar--communication} or \emph{dual-function radar--communication} systems, which emphasize dual-functional waveforms and link-level co-design; and third, related terms such as Joint Communication and Sensing (JCAS), Joint Radar and Communication (JRC), and Radar--Communication (RadCom), which are close to ISAC in meaning but differ slightly in emphasis.

\paragraph{1) Difference between ISAC and Radar--Communication Coexistence}

\begin{sloppypar}
Radar--Communication Coexistence (RCC) is one of the earliest and most widely studied directions in the interaction between radar and communication systems. Its core question is the following: when communication and radar systems, originally designed independently, are forced to share the same or adjacent spectrum bands because of spectrum scarcity or deployment proximity, how can interference modeling, resource coordination, and signal processing be used to enable them to operate jointly with acceptable performance degradation \cite{Chiriyath2017RadarCommCoexistence,Liu2020JointRadarComm}? Typical research topics include analyzing and suppressing radar interference on communication receivers, evaluating the impact of communication signals on radar detection and parameter estimation, and designing coexistence strategies based on power control, beam nulling, cognitive spectrum access, frequency planning, and space--time interference suppression \cite{Deng2013InterferenceMitigation,Sodagari2012ProjectionBasedCoexistence,Zheng2019RadarCommunicationCoexistence}.
\end{sloppypar}

In essence, RCC deals with the \textbf{compatibility problem between two independent systems}. Under this framework, communication and radar systems each serve their own separate objectives, and their main relationship is to avoid damaging each other's operation or to coordinate under conflict. Thus, the starting point of coexistence research is conflict management, and its main goal is to reduce mutual interference so that both systems can function as normally as possible. By contrast, ISAC is concerned with \textbf{simultaneously realizing communication and sensing functionalities within a unified system framework}. Its goal is not merely to reduce friction between two systems, but to achieve joint design, coordinated optimization, and performance enhancement under a common system objective \cite{Liu2022ISAC}. In short, coexistence emphasizes ``how two systems can peacefully coexist,'' whereas ISAC emphasizes ``how one system can natively possess both capabilities.'' This distinction is the primary prerequisite for understanding the conceptual boundary of ISAC.

\paragraph{2) Relationship between ISAC and Dual-Function Radar--Communication / Joint Radar and Communication}

Dual-Function Radar--Communication (DFRC) systems are generally regarded as one of the most important technical origins of ISAC and one of the most representative branches of communication--sensing fusion research \cite{Hassanien2015dual,Liu2020JointRadarComm}. The core idea of DFRC is to design a dual-functional signal or transmission mechanism that can simultaneously support information transmission and target sensing, so that the same transmitted waveform can be demodulated by a communication receiver while also enabling a radar receiver to extract target range, velocity, angle, and other parameters from the echoes. Based on this idea, researchers have proposed many representative physical-layer methods, such as OFDM-based dual-function waveform design, joint beamforming, embedded information modulation, symbol-level waveform design, and Multiple-Input Multiple-Output (MIMO) radar--communication integrated architectures \cite{Hassanien2015dual,Zhang2021ISAC6G}.

Conceptually, DFRC may be understood as an important realization of ISAC.
At a technical level, it mainly concerns physical-layer signal design and
link-level transceiver design.
Many problems studied in DFRC, such as dual-functional waveform design,
joint transmission strategies, and communication--radar performance tradeoffs,
constitute key components of the ISAC technical framework.
However, the conceptual scope of ISAC is usually broader than that of DFRC
\cite{Liu2022ISAC,Gonzalez2024integrated}.
DFRC tends to focus more specifically on how a single signal or a single
transmission platform can simultaneously serve radar and communication purposes.
Its emphasis is mainly on the physical layer, waveform layer, array processing
layer, and link-level performance.
By contrast, ISAC not only includes these issues, but also emphasizes
higher-layer problems, such as network architecture, protocol mechanisms,
resource scheduling, multi-node cooperative sensing, sensing information
feedback and fusion, and service-loop construction.
Therefore, DFRC can be considered an important subset or implementation
direction of ISAC, but the two are not equivalent.
The former is more link-level and waveform-oriented, whereas the latter
represents a broader system-level framework.

It should further be noted that the term JRC overlaps substantially with DFRC in the literature, and both are often used to describe the joint design of radar and communication functions \cite{Liu2020JointRadarComm,Zhang2021ISAC6G}. In many cases, JRC emphasizes the idea of \emph{joint design}, whereas DFRC emphasizes the implementation feature that a single platform or waveform performs dual functions. However, this distinction is not absolute, and terminology varies across the literature. Therefore, when reading related work, one should interpret these terms in context rather than making rigid distinctions based solely on the literal wording.

\paragraph{3) Differences among ISAC, JCAS, JRC, RadCom, and related terms}
\begin{sloppypar}
In addition to DFRC, several other terms with meanings close to ISAC frequently appear in the literature, such as JCAS, JRC, and RadCom \cite{Liu2022ISAC,Gonzalez2024integrated}. These terms overlap to some degree, but they also reflect differences in research emphasis and disciplinary tradition.
\end{sloppypar}

Among them, JCAS is semantically the closest to ISAC, and in many papers the two are used almost interchangeably. Relatively speaking, JCAS emphasizes the word ``joint,'' highlighting the joint nature of communication and sensing design, whereas the word ``integrated'' in ISAC emphasizes deeper fusion at the levels of system architecture, resource management, protocol mechanisms, and service workflows. Thus, in terms of subtle semantic difference, ISAC typically implies stronger system-level integration than JCAS, although in practice the two often refer to closely related or overlapping research scopes.

JRC, by contrast, explicitly uses ``radar'' as the representative form of sensing, and is therefore more common in traditional radar--communication fusion research, especially when the main tasks involve radar detection, parameter estimation, and target tracking \cite{Liu2020JointRadarComm,Hassanien2015dual}. ISAC, by comparison, employs the broader term ``sensing,'' which may encompass not only radar targets in the conventional sense but also users, devices, environments, propagation structures, and even human activities. Therefore, JRC can often be regarded conceptually as a more specific, radar-oriented subset of ISAC.

RadCom or Radar--Communication is usually a broad field-level label that can refer in general to research on the integration of radar and communications, but its specific meaning depends heavily on context. In some papers it refers to coexistence problems, while in others it refers to joint design. Its definition is therefore not always used consistently. For this reason, using such terms in a textbook without explicit clarification may easily confuse beginners.

For the above reasons, this textbook consistently adopts ISAC as the core term. This choice is based on three main considerations. First, the word ``integrated'' better captures the deep coupling and holistic design of communication and sensing. Second, the word ``sensing'' is broader than ``radar'' and can cover radar detection, localization, environmental sensing, human activity recognition, and channel-environment sensing, among others. Third, ISAC has become one of the most widely used terms in academia and industry when discussing communication--sensing fusion, and has received sustained attention in 6G vision studies and related standardization discussions \cite{Liu2022ISAC,Gonzalez2024integrated,De2021convergent}.

\paragraph{4) The multilayer meaning of ``Integrated''}

Finally, it should be further emphasized that the word ``integrated'' in ISAC is not a one-dimensional concept, but rather a layered and systematic one. Integration does not simply mean that communication and sensing share a spectrum band or use a common hardware platform. Instead, it may appear across multiple layers ranging from the physical layer to the service layer \cite{Liu2022ISAC,Zhang2021ISAC6G,Liu2020JointRadarComm}. From the perspective of teaching and system analysis, this integration may be roughly summarized into the following levels:

\begin{itemize}
    \item \textbf{Signal-level integration:} communication and sensing use a unified waveform, or waveform design is jointly optimized to simultaneously support information transmission and parameter sensing.

    \item \textbf{Hardware-level integration:} antenna arrays, RF chains, ADCs/DACs, and baseband processing platforms are jointly used by communication and sensing, and the hardware architecture is designed to satisfy the requirements of both functions.

    \item \textbf{Algorithm-level integration:} signal processing and optimization algorithms simultaneously serve both communication and sensing objectives, such as joint channel estimation and target detection, joint beamforming, and joint resource scheduling with state sensing.

    \item \textbf{Network-level integration:} sensing is incorporated into the overall network architecture, including sensing task triggering and scheduling, sensing data uploading and distribution, multi-node cooperative sensing, and feedback from sensing results to communication network optimization.
    
    \begin{sloppypar}
    \item \textbf{Service-level integration:} communication and sensing services are jointly orchestrated for upper-layer tasks, forming a closed-loop service workflow of ``sensing--communication--computation--decision--control'' to support complex applications such as autonomous driving, industrial control, low-altitude coordination, and smart cities.
    \end{sloppypar}
\end{itemize}

It should be noted that different ISAC systems may exhibit different depths of integration across these levels. Therefore, ISAC in practice is not a binary concept of either integrated or not integrated, but is better understood as a continuous spectrum from shallow to deep fusion. For this reason, judging whether a system belongs to ISAC cannot be based simply on whether it shares part of the spectrum or hardware; rather, one should comprehensively evaluate the degree of integration and the depth of joint design across signals, resources, protocols, networks, and service workflows.

\subsubsection{The Working Definition of ISAC in This Textbook}
\label{subsubsec:ch1-working-definition}

After discussing the general definition of ISAC and its boundaries with related concepts, this textbook further introduces a \emph{working definition} of ISAC for teaching, modeling, and subsequent analysis. Such a working definition is necessary because ISAC is still a rapidly evolving research area, and its concrete realizations, system boundaries, and application emphases are not completely consistent across the literature. For a textbook, a general definition alone is insufficient; it is also necessary to explicitly state the basic standpoint, abstraction level, and analytical scope adopted throughout the book, so that the subsequent chapters can remain consistent in terms of theoretical models, system architectures, and performance metrics.

Based on this consideration, this textbook adopts the following interpretive framework for ISAC.
\begin{enumerate}[label=(\arabic*)]
    \item \textbf{ISAC is regarded as one of the key native capabilities of 6G and future wireless networks:}
    
    This textbook discusses ISAC within the broader technical framework of 6G and future wireless networks. Unlike traditional networks, which mainly emphasize connectivity, future wireless networks are widely expected to possess multiple capabilities simultaneously, including communication, sensing, computing, and control, with sensing no longer depending solely on external standalone sensor systems but gradually becoming one of the native capabilities of network nodes and infrastructure \cite{Liu2022ISAC,Gonzalez2024integrated,Saad2020Vision6G}. Under this positioning, the discussion of ISAC in this textbook is not limited to physical-layer signal design for point-to-point links, but also covers link-layer, network-layer, and system-level architectural issues, including cellular networking, multi-node cooperative sensing, edge computing support, and device--edge--cloud collaborative processing. In other words, the ISAC considered in this textbook is not merely a narrow ``dual-function waveform design problem,'' but a system-level capability framework for the evolution of future wireless networks.
    
    \item \textbf{The focus is on system-level problems of coordinated communication and sensing design:}
    
    The core concern of this textbook is not to discuss a single waveform, a single algorithm, or a single application in isolation, but to understand how communication and sensing functions can be jointly designed within a unified system so as to optimize overall performance and enhance task support capability \cite{Liu2020JointRadarComm,Zhang2021ISAC6G}. Specifically, the subsequent chapters will revolve around the following questions:
    \begin{itemize}
        \item How can signal models and transmit waveforms be designed to simultaneously support communication and sensing?
        \item How can wireless resources in time, frequency, space, and power be shared, allocated, and jointly optimized between communication and sensing?
        \item How can joint beamforming strategies be designed in MIMO, multi-antenna, and large-array systems to balance communication quality and sensing performance?
        \item How can a unified performance analysis framework be established to characterize the intrinsic tradeoffs among communication rate, bit error rate, latency, sensing detection probability, estimation accuracy, and resolution?
        \item How can network architectures support large-scale, multi-node, distributed, and cooperative ISAC services?
    \end{itemize}
    These questions correspond, respectively, to the signal layer, resource layer, array processing layer, performance evaluation layer, and network architecture layer in ISAC research, and together they form the main technical thread of this textbook.
    
    \item \textbf{No single optimal waveform, architecture, or application scenario is presupposed:}
    
    ISAC is an open and diverse technological framework, and there is no universally optimal waveform, universally optimal architecture, or single standard application paradigm \cite{Liu2022ISAC,Liu2020JointRadarComm}. Different frequency bands, node capabilities, network deployment modes, and service objectives all lead to substantially different ISAC system designs. Therefore, this textbook does not reduce ISAC to any fixed implementation, but instead introduces multiple representative technical pathways based on common principles, including but not limited to OFDM-based ISAC waveform design, MIMO-based joint beamforming schemes, sensing networking and cooperative processing in cellular systems, high-resolution sensing-and-communication systems in the mmWave/THz bands, and ISAC application frameworks for vehicular networks, the low-altitude economy, smart factories, and intelligent spaces \cite{Liu2022ISAC,Gonzalez2024integrated}. In this way, readers can understand both the common principles of ISAC and its differentiated implementation forms under different technical conditions and application scenarios.
\end{enumerate}

This working framework not only defines the scope of the textbook, but also lays the foundation for a unified mathematical modeling and performance analysis language in later chapters. Specifically, in the subsequent theoretical analyses of this book, communication performance will mainly be described by metrics such as channel capacity, achievable rate, bit error rate, reliability, and transmission latency; sensing performance will mainly be characterized by metrics such as detection probability, false alarm probability, range/velocity/angle resolution, estimation error, and the CRLB \cite{Liu2022ISAC,Zhang2021ISAC6G}. When communication and sensing need to operate jointly in a unified system, the two classes of metrics generally do not exhibit a simple monotonic relationship. Instead, their tradeoffs and synergy boundaries must often be characterized using tools such as multiobjective optimization, constrained optimization, and Pareto optimality \cite{Liu2020JointRadarComm}. In terms of system modeling, this textbook will strive to use a unified mathematical language to describe the transmit signal model, propagation channel model, target/scatterer model, observation model, and receiver processing chain, and on this basis gradually establish a joint performance evaluation framework for analyzing the relationships among rate, reliability, latency, localization accuracy, detection probability, estimation error, and energy consumption.

It should be emphasized that this working definition does not restrict ISAC to any narrow implementation. Rather, it means that in the subsequent chapters, any system that satisfies the following essential characteristic will be considered within the scope of ISAC: \emph{under a unified system framework, communication and sensing functions provide joint information transmission and environment/state sensing services through shared or jointly designed signals, resources, hardware, protocols, and network mechanisms, while exhibiting bidirectional empowerment between the two} \cite{Liu2022ISAC,Gonzalez2024integrated}. Conversely, schemes that only realize physical co-location, simple spectrum coexistence, or loose functional concatenation without a unified system objective and joint design mechanism will generally not be regarded in this textbook as part of the strict core category of ISAC.

Based on the above conventions, this textbook adopts the following working definition:
\begin{quote}
\textbf{ISAC is a key native capability for 6G and future wireless networks. It refers to the joint provisioning of information transmission and environment/state sensing services under a unified system framework through shared or jointly designed signals, resources, hardware, protocols, and network mechanisms, while supporting network intelligence, scene understanding, and service loops through the bidirectional empowerment of communication and sensing.}
\end{quote}

This working definition will serve as the common starting point for the system modeling, algorithm design, performance analysis, and application discussions throughout the subsequent chapters of this book.

\subsection{The Three-Stage Logic of ISAC Evolution}
\label{subsec:ch1-evolution-logic}

The development of ISAC technology did not occur overnight. Rather, it has undergone a gradual evolution from a state in which communication and sensing were independent but forced to coexist, to one in which they are deeply integrated under a unified system framework. To help readers understand this developmental trajectory, this textbook summarizes the evolution of communication--sensing fusion into three stages with clear progressive relationships: the \emph{spectrum coexistence and cross-application stage}, the \emph{dual-functional platform and joint signal design stage}, and the \emph{6G-oriented native integrated sensing and communication stage}. These three stages should not be understood as strictly isolated historical episodes, but rather as a technological spectrum in which the focus of research deepens and the degree of system integration continuously increases. Understanding the core objectives, representative solutions, and inherent limitations of each stage helps reveal the internal logic by which ISAC evolved from ``coexistence'' to ``integration.''

\subsubsection{Stage 1: Spectrum Coexistence and Cross-Application Phase}
\label{subsubsec:ch1-stage1}

The first stage in the evolution toward ISAC may be summarized as the \emph{spectrum coexistence and cross-application phase}. This stage roughly originated when communication systems and radar systems began to experience significant spectrum conflicts. The core question in this phase was not ``how to build one system that performs both communication and sensing,'' but rather ``when two originally independent systems must share the same spectrum or occupy adjacent bands, how can mutual interference be minimized so that both systems can continue to function properly?'' \cite{Liu2020JointRadarComm,Chiriyath2017RadarCommCoexistence}. Hence, the defining characteristic of this stage is that communication and radar systems remain independent in terms of objectives, architectures, waveforms, protocols, and performance evaluation, while some degree of coordination and compatibility is needed at the spectrum level.

From a system perspective, the central contradiction in this stage arises from the tension between spectrum scarcity and the independent deployment of two separate systems. Communication and radar each use their own signal formats, hardware platforms, receiver processing chains, and optimization objectives: the former usually aims for higher transmission rate, lower bit error rate, and better link reliability, whereas the latter focuses on detection probability, false alarm probability, parameter estimation accuracy, and spatial resolution. Under such circumstances, research and engineering practice were mainly concerned with building \emph{compatibility mechanisms}, rather than building a unified system \cite{ITU_R_RR2024}. In other words, the ``jointness'' at this stage exists only in spectrum coordination, not yet in integrated signal design, resource optimization, or network architecture.

Typical research problems in this stage include the following:
\begin{itemize}
    \item \textbf{Spectrum conflict analysis:} quantitatively evaluating the mutual impact between communication and radar systems under co-channel or adjacent-channel operation, including statistical modeling of interference signals, interference power calculation, received SINR analysis, and evaluation of the effects of interference on communication bit error rate, radar detection performance, and parameter estimation accuracy \cite{Chiriyath2017RadarCommCoexistence}.
    
    \item \textbf{Interference management:} designing mechanisms to reduce mutual interference, such as transmit power control, directional beam isolation, guard band design, dynamic spectrum coordination, space--time filtering, and interference cancellation at the receiver \cite{Liu2020JointRadarComm}.
    
    \item \textbf{Coexistence mechanism design:} establishing operating rules for communication and radar systems under shared spectrum, including sensing- and cognition-based dynamic spectrum access protocols, spectrum occupancy detection, shared timing coordination, and spectrum reuse strategies that satisfy regulatory requirements \cite{Chiriyath2017RadarCommCoexistence}.
\end{itemize}

In essence, during the coexistence stage, communication and radar still follow the philosophy of ``independent optimization with compatible operation.'' Communication systems are designed around transmission efficiency and link quality, whereas radar systems are designed around detection and estimation performance. There is no true unified objective function between them, nor is there a joint optimization mechanism spanning transmission, reception, resource scheduling, and protocol control. Thus, the technical boundary of this stage can be summarized as follows: \textbf{coordination exists at the level of interference management, but true integration at the level of system design is absent}. This is also the fundamental difference between this stage and the later stages of ISAC.

Several representative technical solutions appeared in this coexistence stage. Although they did not yet achieve integrated design, they laid the problem and methodological foundations for later deep fusion of communication and sensing.
\begin{itemize}
    \item \textbf{Time-division sharing:} communication and radar systems alternately use the same frequency band in the time domain, for example by scheduling communication transmission during radar pulse intervals or assigning predefined time slots separately to sensing and communication tasks. This approach is relatively simple to implement, offers strong isolation between systems, and is convenient for regulatory compliance and deployment. However, its drawback is that the spectrum cannot be truly used simultaneously, and the improvement in overall spectral efficiency is limited \cite{Chiriyath2017RadarCommCoexistence}.
    
    \item \textbf{Frequency-division isolation:} within a shared or adjacent spectrum range, communication and radar are allocated different subbands, and adjacent-channel interference is reduced through guard bands, filtering, and frequency planning. This approach is straightforward in engineering practice and is compatible with existing spectrum management frameworks. However, it is still essentially resource partitioning rather than functional integration, and therefore cannot fully exploit the potential of wideband spectrum resources \cite{ITU_R_RR2024}.
    
    \item \textbf{Coexistence-aware sensing or communication processing:} the receiver or signal processing stage explicitly takes into account the interference structure generated by the other system. For example, a radar receiver may model communication signals as structured interference and suppress them, or a communication receiver may detect, estimate, and cancel radar pulse interference. Such methods improve coexistence capability without changing the independence of the two systems \cite{Chiriyath2017RadarCommCoexistence,Liu2020JointRadarComm}.
    
    \item \textbf{Opportunistic spectrum reuse:} inspired by cognitive radio, the system performs spectrum sensing, occupancy detection, or environmental cognition to determine whether a frequency band is idle or under low interference, and then opportunistically reuses that band under certain constraints. In theory, this approach can improve spectrum utilization, but it places stringent requirements on sensing accuracy, fast access mechanisms, and regulatory compatibility \cite{Chiriyath2017RadarCommCoexistence,Haykin2005CognitiveRadio}.
\end{itemize}

Although these coexistence solutions alleviate spectrum conflict to some extent, their limitations are also evident. First, under the coexistence framework, improving the performance of one system often comes at the expense of degrading the performance of the other, leaving limited room for overall performance improvement. Second, since communication and radar systems remain independent in terms of hardware, waveforms, protocols, and information processing chains, they are unable to form a unified architecture or realize deep resource sharing. Third, the lack of substantive bidirectional information exchange and joint optimization prevents the release of the potential synergistic gains of sensing-assisted communication and communication-empowered sensing. Finally, from the engineering perspective, maintaining two separate systems together with their coordination mechanisms also introduces additional complexity and deployment cost \cite{Liu2020JointRadarComm,Liu2022ISAC}. These inherent limitations drove research beyond the question of ``how to coexist'' toward the question of ``how to realize dual functionality on the same platform and within the same signal,'' thereby leading to the second stage of ISAC evolution.

\subsubsection{Stage 2: Dual-Functional Platform and Joint Signal Design Phase}
\label{subsubsec:ch1-stage2}

If the keyword of the first stage is ``coexistence,'' then the keyword of the second stage is ``joint design.'' This stage marks the first substantial leap from compatibility between two independent systems to the realization of dual functionality on the same platform or within the same signal. It therefore constitutes the first real transition from coexistence to integration in the evolution toward ISAC. The most representative idea of this stage can be summarized as ``one signal, dual purposes'' or ``one platform, dual functions'': communication and sensing no longer rely on two separate transmitters, receivers, and waveform systems, but begin to share the same transmit waveform, antenna array, RF front end, and even parts of the receiver processing chain.

From the perspective of technical objectives, the core issue of this stage is no longer simply reducing interference between two systems, but rather how to jointly design a single signal or platform so that it simultaneously satisfies the performance requirements of both communication and sensing. It is in this stage that DFRC, JRC, and early physical-layer ISAC methods gradually developed into a relatively complete technical framework \cite{Liu2020JointRadarComm,Zhang2021ISAC6G,Hassanien2015dual}. Representative research directions in this stage mainly include the following:
\begin{itemize}
    \item \textbf{Dual-function waveform design:} how to design a transmit waveform that is simultaneously suitable for data transmission and target sensing, so that it has sufficient information-carrying capability while also exhibiting proper time--frequency characteristics for target detection, parameter estimation, and range--velocity resolution \cite{Liu2020JointRadarComm,Hassanien2015dual}. The essential challenge here is that communication and sensing often impose different waveform requirements: communication prefers strong modulation flexibility, high spectral efficiency, and robustness against multipath, whereas sensing is more concerned with autocorrelation properties, ambiguity function shape, bandwidth utilization, and target parameter identifiability. Thus, dual-function waveform design is fundamentally a joint optimization problem seeking a balance between communication achievable rate and sensing estimation performance.
    
    \item \textbf{Joint beamforming:} in multi-antenna and MIMO systems, how to design the transmit precoding matrix or beamforming vector so that it can simultaneously provide communication users with sufficient Signal-to-Interference-plus-Noise Ratio (SINR) or rate guarantees while also forming the desired illumination power, spatial focusing capability, or angular resolution in the target direction \cite{Liu2020JointRadarComm,Liu2022ISAC}. This problem is inherently multiobjective, because communication beams often favor user directions, whereas sensing beams may need to cover target regions or create specific beampattern structures, leading to conflicts in the allocation of spatial resources.
    
    \item \textbf{Shared resource scheduling and joint optimization:} when communication and sensing share the same set of time, frequency, space, and power resources, how to optimally allocate these resources under system constraints becomes another core issue of this stage \cite{Zhang2021ISAC6G,Liu2022ISAC}. Such problems may arise in frame-level pilot/data/sensing allocation, subcarrier selection, power loading, array degree-of-freedom allocation, and time-slot structure configuration. In essence, these problems require a unified framework for analyzing tradeoffs between communication and sensing performance metrics.
\end{itemize}

Although the second stage already achieved substantial joint design of communication and sensing at the physical and link layers, its research boundary was still largely limited to single-link, single-node, or single-cell scenarios. Specifically, joint waveform design was usually studied for point-to-point or point-to-multipoint links, joint beamforming was mostly analyzed within a single base station or transmission platform, and performance evaluation was mainly based on link-level theoretical analysis or simulation \cite{Liu2020JointRadarComm,Liu2022ISAC}. By contrast, issues such as multicell cooperative sensing, network-level sensing task scheduling, sensing result aggregation and fusion, end-to-end orchestration of sensing services, and sensing--edge-computing collaboration had not yet become the main focus of this stage. Therefore, from the perspective of technical evolution, although the second stage laid the core theoretical foundations of ISAC, its system perspective remained primarily link-level and platform-level rather than network-native.

In terms of implementation, several representative types of solutions emerged in this stage. These solutions constitute the most important technical foundation for later ISAC research.
\begin{itemize}
    \item \textbf{OFDM-based ISAC schemes:} OFDM is one of the most representative dual-functional signal carriers in this stage because it naturally possesses a regular time--frequency structure, mature communication transceiver mechanisms, and convenient processing properties for parameter estimation \cite{Liu2022ISAC,Sturm2011WaveformDesign,Zhou2022integrated}. In such schemes, the OFDM waveform carries communication data while target range and velocity are extracted from the received echoes through joint processing across the subcarrier and symbol dimensions. Since the two-dimensional time--frequency structure of OFDM corresponds naturally to the range--Doppler processing in radar, many methods can achieve efficient parameter estimation using Two-Dimensional Fast Fourier Transform (2D-FFT) or its variants. The significance of these schemes lies in the direct demonstration that mainstream communication waveforms can support basic sensing tasks with relatively minor modifications, thereby strongly promoting the evolution of communication systems into ISAC platforms.
    
    \item \textbf{DFRC waveform design:} another important technical route is to embed communication information into traditional radar waveforms, or to explicitly impose radar-oriented structural constraints on communication signals, thereby forming truly dual-functional waveforms \cite{Hassanien2015dual,Blunt2016overview,Blunt2010Embedding}. Such methods include, but are not limited to, embedding modulation information into pulse radar or LFM signals, realizing spatial modulation by antenna indexing or beam switching, optimizing the ambiguity function of the transmit waveform under communication quality constraints, and jointly designing information embedding mechanisms within MIMO radar waveforms. Compared with directly reusing existing communication waveforms, DFRC methods often emphasize constructing the dual-functional signal from the radar-performance perspective, which can offer advantages in target detection and angle control, but may face challenges in communication compatibility, standards adaptation, and system complexity.
    
    \item \textbf{MIMO joint design and beamforming:} with the development of multi-antenna technology, MIMO became one of the most important platforms for joint design in this stage \cite{Liu2020JointRadarComm,Li2023integrated}. In such schemes, the transmitter jointly designs the precoding matrix, covariance matrix, or beamforming vector so that user SINR, rate, or bit error rate constraints are simultaneously satisfied together with the sensing-side requirements on beampattern shape, illumination power, and angular discriminability. These problems are often formulated as multiobjective optimization, constrained optimization, or robust optimization problems, and are commonly solved using Semidefinite Relaxation, Alternating Optimization (AO), weighted-sum methods, penalty methods, and block coordinate descent. The importance of MIMO joint design lies in the fact that it translates the conflict and synergy between communication and sensing into a unified optimization framework, thereby laying the methodological foundation for later developments in massive MIMO, mmWave beam management, and network-level cooperative sensing.
   
    \begin{sloppypar}
    \item \textbf{Joint performance optimization for link-level tradeoffs:} beyond specific waveform and array designs, the second stage also gradually established the theme of \emph{communication--sensing performance tradeoff} as a central research topic \cite{Liu2020JointRadarComm,Liu2022ISAC}. Representative works include deriving tradeoff relationships between communication rate and detection performance, characterizing constraint relationships between bit error rate and radar parameter estimation accuracy, and analyzing the achievable performance region of dual-function systems through Pareto boundaries, multiobjective optimization, and constrained optimization. It is within this class of studies that ISAC began to form its own distinctive language and methodology of performance analysis.
    \end{sloppypar}
\end{itemize}

Overall, the significance of the second stage is that it systematically demonstrated, for the first time at the physical and link layers, that communication and sensing need not merely coexist, but can be genuinely fused on the same signal, the same platform, and within the same optimization framework \cite{Liu2020JointRadarComm,Liu2022ISAC,Hassanien2015dual}. Through this stage, researchers gradually developed a set of core theories, mathematical models, and algorithmic tools centered on dual-functional waveforms, joint beamforming, joint resource optimization, and communication--sensing performance tradeoffs, and verified the feasibility and effectiveness of mainstream communication technologies such as OFDM and MIMO as carriers of ISAC. However, this stage remained mainly focused on link-level and platform-level problems, with limited attention to network-level cooperation, service-oriented sensing, and native architectural support. Therefore, with the emergence of the 6G vision and the demand for network intelligence, ISAC naturally evolved toward the third stage, namely the stage of native integrated sensing and communication as a built-in network capability.

\subsubsection{Stage 3: 6G Native Integrated Sensing and Communication Phase}
\label{subsubsec:ch1-stage3}

The third stage represents the frontier direction of ISAC development and is widely regarded as an important vision for 6G wireless networks \cite{Liu2022ISAC,Gonzalez2024integrated,Saad2020Vision6G}. The most fundamental characteristic of this stage is that sensing is no longer treated as an external add-on to communication systems, but is elevated to a \emph{native capability} of wireless networks. The term ``native'' means that sensing is incorporated into the overall network design from the outset and is natively supported at multiple levels, including standards and protocols, network architecture, resource scheduling, service orchestration, and end-to-end application support. Compared with the first two stages, ISAC in this stage no longer remains mainly at the level of physical-layer waveform design and link-level optimization, but expands into a system-level design problem covering the physical layer, link layer, network layer, and even the application and service layers \cite{Liu2022ISAC,Zhang2020perceptive}.

From the technical standpoint, the essence of this evolution is the expansion of the ISAC research object from a \emph{single-link dual-functional platform} to a \emph{networked, service-oriented, and intelligent communication--sensing system}. In other words, while the second stage asks ``how can one platform communicate and sense simultaneously?'' the third stage further asks ``how can one network natively possess sensing capability and provide it as a service to support complex applications?'' \cite{Zhang2020perceptive,Cui2024integrated}. Under this view, the integration boundary between communication and sensing is significantly broadened, and network-level cooperation, computing collaboration, service loops, and intelligent control become essential parts of ISAC.

Representative technical problems in the third stage include the following:
\begin{itemize}
    \item \textbf{Network-level cooperative sensing:} how can observation resources from multiple base stations, access points, or sensing nodes in a cellular network be coordinated to achieve wider-area, more robust, and higher-resolution cooperative sensing \cite{Zhang2020perceptive,Cui2024integrated}? In this context, issues such as inter-node time synchronization, frequency synchronization, phase consistency, observation data feedback, distributed fusion, and joint target tracking all significantly affect system performance.
    
    \item \textbf{Cellular ISAC networking:} how can sensing functionality be deployed and managed within a cellular network topology so that sensing coverage and sensing resolution can be coordinated with communication coverage, user capacity, and interference management \cite{Liu2022ISAC,Zhang2020perceptive}? This involves many aspects, including the relationship between sensing cells and communication cells, sensing task triggering mechanisms, air-interface resource allocation, and functional splitting at the radio access network side.
    
    \item \textbf{Device--edge--cloud collaborative processing:} because raw sensing data are often large in scale and require timely processing, how to distribute sensing data processing, compression, fusion, and decision tasks among devices, edge nodes, and cloud platforms becomes an important system-level problem in this stage \cite{Gonzalez2024integrated,Saad2020Vision6G}. This issue affects not only sensing accuracy but also backhaul load, system latency, and energy consumption.
    
    \item \textbf{Service-loop support:} how to feed sensing results back in real time to communication, computing, and control systems so as to form a closed-loop service chain of ``sensing--transmission--computation--decision--control'' is one of the key features distinguishing this stage from the previous two \cite{Liu2022ISAC,Saad2020Vision6G}. In scenarios such as autonomous driving, industrial control, unmanned system collaboration, and digital twins, sensing is no longer merely an independent output, but the initial step and critical input of the entire task loop.
\end{itemize}

\begin{sloppypar}
Compared with the previous two stages, the technical boundary of the third stage has clearly expanded from link-level fusion to system-level, network-level, and service-level fusion. Specifically, the design object of ISAC is no longer confined to a single base station, a single platform, or a single link, but extends to the joint optimization of the entire network system. This expansion manifests itself at least in three respects \cite{Zhang2020perceptive,Cui2024integrated}. First, at the resource management level, sensing resource scheduling, multicell interference and coordination, sensing coverage planning, and cross-node observation cooperation must all be handled at the network-wide scale. Second, at the architecture and protocol levels, sensing must be incorporated into network functional entities and control-plane design, for example by introducing sensing task management, sensing data fusion and analysis, and service orchestration modules at the core network or edge side, while natively supporting sensing signal transmission, observation, and feedback at the air interface and control protocols of the access network. Third, at the service level, communication services and sensing services no longer run along two independent pipelines, but are jointly orchestrated and delivered for upper-layer tasks, allowing applications to request combined ``connectivity + sensing + computing'' capabilities on demand.
\end{sloppypar}

At this stage, ISAC systems begin to exhibit new capabilities and system-level value that were difficult to achieve in the first two stages.
\begin{itemize}
    \item \textbf{Sensing-assisted communication:} sensing results are used in reverse to improve communication performance. For example, by obtaining user location, movement trajectories, environmental blockage states, and spatial propagation structures, the network can perform beam prediction, handover preparation, blockage warning, and resource scheduling in advance, thereby reducing beam training overhead, improving link stability, and lowering handover latency \cite{Liu2022ISAC,Zhang2020perceptive}. This capability is particularly important in the mmWave and THz bands, where narrow-beam high-frequency communication is highly sensitive to spatial state information.
    
    \item \textbf{Communication-empowered sensing:} the existing infrastructure, signal resources, and spatiotemporal coverage of communication networks are used to realize large-scale and continuous environmental sensing \cite{Liu2022ISAC,Gonzalez2024integrated}. For example, the multi-base-station deployment of cellular networks naturally provides a distributed observation network for cooperative sensing, continuously available downlink and uplink communication signals can serve as opportunistic illumination sources, and network synchronization and backhaul mechanisms provide the basis for multi-node data fusion. Compared with traditional standalone sensing systems, this approach is more conducive to wide-area coverage and infrastructure reuse.
    
    \item \textbf{Sensing-driven network intelligence:} physical-world information obtained from sensing is incorporated into intelligent network decision-making, such that resource scheduling, mobility management, load balancing, beam control, and quality-of-service assurance no longer rely solely on communication-side feedback, but can instead be optimized based on real-time understanding of the physical environment \cite{Saad2020Vision6G,Cui2024integrated}. This also creates a much tighter coupling between Artificial Intelligence / Machine Learning (AI/ML) techniques and ISAC, pushing the network from an ``adaptive connectivity system'' toward an ``intelligent system with environmental awareness.''
    \begin{sloppypar}
    \item \textbf{Service loop:} a complete closed loop of ``sensing--communication--computation--decision--control--re-sensing'' is formed \cite{Saad2020Vision6G,Liu2022ISAC}. For instance, in industrial automation, when the network-side sensing system detects abnormal equipment operation, the observed information can be quickly transmitted to an edge computing node through the communication link, where condition assessment and control decisions are made; the control commands are then delivered to the actuators, and the sensing function continues to monitor the outcome and provide feedback. Such a closed-loop capability means that ISAC is no longer merely a functional extension of communication systems, but becomes an important foundational capability for digitalized, intelligent, and automated applications.
    \end{sloppypar}
\end{itemize}

From the perspective of technological evolution, the third stage marks a fundamental transition of ISAC from a \emph{link-level dual-functional technology module} to a \emph{network-level native capability framework} \cite{Liu2022ISAC,Zhang2020perceptive,Cui2024integrated}. In this stage, ISAC is no longer merely a signal processing problem in physical-layer research, but also becomes a topic of common concern in network architecture design, protocol evolution, edge intelligence, standardization, and industrial ecosystem development. At the same time, this stage raises many open questions that still require in-depth study, such as the fundamental performance limits of network-level ISAC, optimal task allocation mechanisms for multi-node cooperative sensing, security and privacy protection of sensing data, and business models and resource pricing methods for sensing-as-a-service \cite{Gonzalez2024integrated,Cui2024integrated}.

From the standpoint of standardization and industrial evolution, ISAC is also moving gradually from academic research toward engineering system construction. In recent years, the Third Generation Partnership Project (3GPP) has launched discussions on ISAC-related capabilities and initiated related study work in Release 19. At the same time, in International Mobile Telecommunications 2030 (IMT-2030)-oriented 6G vision studies, the integration of sensing and communications has already been identified as an important direction for future wireless networks \cite{IMT2030Framework,ThreeGPPRel19ISAC}. These developments indicate that the third stage is not merely an abstract vision, but is progressively entering a phase where it can be researched, implemented, and promoted through standardization.

\subsection{The Core Driving Forces Behind the Three-Stage Evolution}
\label{subsec:ch1-evolution-driving-force}

The evolution of ISAC along the three stages described above is not the result of several isolated technical solutions being accidentally assembled together. Rather, it has been jointly driven by a series of long-standing and increasingly strong structural forces. These driving forces include objective constraints at the spectrum and hardware-resource levels, as well as capability demands imposed by future network positioning, emerging application scenarios, and the trend toward intelligent systems. Understanding these drivers helps reveal, at a deeper level, the internal logic by which ISAC evolved from coexistence to integration, and further from link-level integration to a network-level native capability.

\subsubsection{Spectral Efficiency and Resource Reuse Requirements}
\label{subsubsec:ch1-driving-force-spectrum}

Spectrum scarcity is one of the most direct drivers of ISAC development. As global mobile data traffic continues to increase, the demand for wireless spectrum keeps rising, while high-quality contiguous spectrum suitable for mobile communications, especially in the sub-6 GHz range, has become increasingly scarce \cite{ITU_R_RR2024,miit_radio_freq_2026_en,Saad2020Vision6G}. At the same time, in order to support higher throughput and larger system capacity, communication systems are moving toward the mmWave and even THz bands, where there is growing overlap or adjacency with the operating frequencies of radar, remote sensing, and radiolocation systems \cite{Liu2022ISAC}. It is under this background that spectrum conflict between communication and sensing first drove the development of spectrum coexistence and interference management technologies in the first stage.

However, coexistence can only alleviate conflict to a limited degree and cannot fundamentally improve spectrum utilization efficiency. The basic reason is that, under the coexistence framework, communication and sensing systems still remain competitors in essence: improving the performance of one system often compresses the usable resources or raises the interference level for the other \cite{Liu2020JointRadarComm,Chiriyath2017RadarCommCoexistence}. For this reason, simple coexistence mechanisms cannot fully exploit the potential of shared spectrum. By contrast, if communication and sensing can share the same transmit signal, the same time--frequency resources, and even the same hardware platform, then the same physical resources can simultaneously serve the dual purposes of information transmission and environmental sensing, yielding resource reuse gains beyond those achievable by simple time-sharing or frequency-sharing. This transition from ``spectrum sharing'' to ``signal and function sharing'' constitutes one of the most fundamental motivations for the joint signal design of the second stage \cite{Liu2022ISAC,Hassanien2015dual}.

At a deeper level, as 6G networks are expected to simultaneously support communication, sensing, localization, computing, and control, the traditional system design paradigm of ``one function occupying one dedicated set of resources'' is becoming unsustainable \cite{Gonzalez2024integrated,Saad2020Vision6G}. Future networks must allow time, frequency, space, and power resources to support multiple functions at once. This means that system design philosophy must evolve from ``function-specific resource allocation'' toward ``resource reuse, functional coordination, and unified scheduling.'' ISAC is precisely the representative manifestation of this trend in the wireless air interface and network architecture.

\subsubsection{Hardware Integration and the Development of Reconfigurable Platforms}
\label{subsubsec:ch1-driving-force-hardware}

If spectrum pressure provides the external impetus for ISAC, then advances in hardware platforms provide the key physical foundation for its realization. One major reason why communication and sensing have gradually moved from being merely combinable in theory to being implementable in practice is that modern wireless hardware platforms are becoming increasingly reconfigurable, broadband, array-based, and software programmable.

First, the development of Software-Defined Radio (SDR) and general-purpose baseband platforms has made it possible for the same RF and baseband hardware to support different signal schemes and processing chains through software configuration, thereby significantly lowering the barrier to implementing communication and sensing on the same platform \cite{Mitola1995SDR}. This programmability means that it is no longer necessary to deploy completely separate hardware chains for different functions. Instead, multitask reuse can be achieved through software upgrades, parameter reconfiguration, and functional switching. This provides the essential foundation for dual-functional platform design in the second stage.

Second, the development of massive MIMO and phased-array technologies has provided ISAC with abundant spatial degrees of freedom. By configuring a large number of antenna elements, the system can simultaneously form multiple beams with different directions, widths, and power distributions in the spatial domain. This allows high-quality directional coverage for communication users while also providing controllable illumination and angular resolution for sensing targets \cite{Liu2020JointRadarComm,Li2023integrated}. It is precisely this spatial flexibility that makes joint beamforming and spatial resource reuse possible, and that gives practical feasibility to many of the MIMO-based joint design schemes in the second stage.

Third, mmWave and THz communication technologies naturally exhibit physical properties that are highly compatible with sensing requirements. High-frequency systems usually have much larger available bandwidth, which enables higher range resolution. In the classical monostatic setting, for example, the range resolution approximately satisfies
\[
\Delta R = \frac{c}{2B},
\]
where $c$ denotes the speed of light (m/s) and $B$ is the effective signal bandwidth (Hz) \cite{SkolnikRadarBook}. This means that a larger bandwidth yields stronger resolvability between targets in the range dimension. At the same time, mmWave and THz systems usually rely on large antenna arrays and narrow-beam transmission, which are naturally beneficial for improving angular resolution and spatial focusing capability \cite{Rappaport2019WirelessTHz}. Therefore, from the perspective of physical-layer characteristics, high-frequency communication systems inherently possess the potential to perform high-precision sensing tasks. This is also one of the main reasons why ISAC has attracted particular attention in the mmWave and THz bands.

It should also be noted that at lower frequencies, traditional communication equipment and radar equipment often differ significantly in RF architecture, bandwidth requirement, and array structure, making hardware sharing relatively difficult. By contrast, in the mmWave and higher-frequency ranges, communication and sensing systems exhibit increasing convergence in front-end implementation, as both tend to rely on broadband RF chains, phased-array structures, high-speed ADC/DAC modules, and tightly coupled digital beamforming architectures \cite{Liu2022ISAC,Rappaport2019WirelessTHz}. This hardware convergence substantially lowers the engineering barrier to realizing communication and sensing on a common platform, thus pushing ISAC from a conceptually feasible idea toward a practically realizable system.

In addition, the development of new reconfigurable electromagnetic devices has opened up new design space for ISAC. Reconfigurable Intelligent Surfaces (RIS), as a representative example of programmable electromagnetic environment technology, can actively shape the wireless propagation environment by adjusting the amplitude and phase responses of reflecting elements, thereby simultaneously serving communication link enhancement and sensing coverage extension \cite{Chepuri2023integrated}. In ISAC scenarios, RIS can not only improve coverage blind spots, enhance sensing echoes, or create favorable propagation paths, but also provide additional degrees of freedom for the joint optimization of communication and sensing. For this reason, RIS has become one of the important recent research hotspots.

\subsubsection{Network Intelligence and Service-Loop Requirements}
\label{subsubsec:ch1-driving-force-network}

Beyond the driving forces arising from resource and hardware constraints, the changing role of mobile communication networks themselves is also an important factor in the evolution of ISAC. In 5G and earlier generations, the main value of wireless networks lay in providing high-quality connectivity. In the 6G vision, however, networks are expected to evolve further into comprehensive service platforms that integrate connectivity, sensing, computing, control, and intelligence \cite{Gonzalez2024integrated,Saad2020Vision6G}. In this transformation, an increasing number of emerging applications demand closed-loop services of ``sensing--communication--computation--decision--control,'' making sensing rise from a peripheral auxiliary function to a key native capability of the network itself.

For example, in autonomous driving, the system must complete environmental detection, state sharing, path planning, and control execution on a millisecond time scale; in industrial Internet and smart manufacturing scenarios, the system must realize a real-time closed loop of ``monitoring--transmission--analysis--adjustment''; and in remote control and digital operation-and-maintenance scenarios, the network must continuously sense the states of physical devices and environments in order to support highly reliable feedback control \cite{Liu2022ISAC,Saad2020Vision6G}. In such closed loops, sensing is not an optional output, but the starting point and critical input of the entire service flow. If sensing depends entirely on independent systems external to the network, issues such as sensing-data return latency, interface complexity, clock synchronization, and heterogeneous data fusion may become bottlenecks that constrain closed-loop performance. Therefore, embedding sensing into the network architecture becomes an important prerequisite for achieving efficient closed-loop services.

Meanwhile, as AI/ML technologies are increasingly applied to wireless network optimization, network decisions depend more and more on an accurate understanding of environmental state. For example, intelligent beam management requires user location and movement direction information, intelligent resource scheduling requires knowledge of traffic distribution in time and space, and intelligent interference management requires identifying the locations and activity states of potential interference sources \cite{Saad2020Vision6G}. Traditional networks mainly rely on channel feedback, measurement reports, and historical statistics for optimization. Although such information is important, it is often limited in terms of environmental interpretability, real-time responsiveness, and predictive capability. ISAC addresses this limitation by natively introducing environmental sensing into the network, enabling the network to obtain richer and more direct information about the physical world, and thus providing high-value input for network intelligence. This trend directly drives the emergence of \emph{sensing-driven network intelligence} in the third stage.

\subsubsection{6G Vision Requirements for Native Capabilities}
\label{subsubsec:ch1-driving-force-6g-vision}

From an even broader perspective, the 6G vision itself imposes requirements on wireless networks that go beyond traditional connectivity. It is now widely recognized that future networks should not merely provide high-speed transmission and ubiquitous connectivity, but should also possess the ability to sense environments, targets, users, and service states in real time, to understand scenarios, and to support coordinated control \cite{Saad2020Vision6G,IMT2030Framework}. In this context, the shared demands of new scenarios such as ultra-reliable low-latency services, native intelligence, digital twins, the low-altitude economy, and integrated space--air--ground--sea networks are driving communication systems to evolve from possessing a single connectivity capability toward a native capability framework that deeply integrates communication, sensing, computing, and intelligence. ISAC is one of the key technological foundations enabling this evolution \cite{Gonzalez2024integrated,Liu2022ISAC}.

First, ultra-high reliability and low-latency services require stronger prediction and determinism from the network. Traditional wireless networks often rely on reactive adjustment after link quality has already degraded, which makes them insufficient for future services that require millisecond-level or even lower latency and near-zero interruption. Through simultaneous acquisition of user position, movement trajectory, blockage variation, environmental reflector distributions, and potential interference information during communication, ISAC enables the network to perceive channel and topology changes earlier and to perform beam management, handover, and resource allocation proactively \cite{Liu2022ISAC,Zhang2020perceptive}. In this way, the network evolves from merely ``repairing after the fact'' to being able to ``sense in advance, decide in advance, and optimize in advance.''

Second, the concept of native intelligence further strengthens the dependence of networks on environmental information. Intelligent networks must not only transmit data, but also learn, infer, and optimize autonomously based on environmental state and service state. One of the most important categories of input for this purpose is precisely the dynamic sensing information originating from the real physical world \cite{Saad2020Vision6G}. By embedding sensing into network infrastructure, ISAC can continuously acquire multidimensional information about user behavior, scene variation, device status, and spatial distribution, thereby providing training samples and real-time inference input for AI models used in resource scheduling, service prediction, anomaly detection, mobility management, and multitask coordination.

Third, the digital twin vision further amplifies the need for real-time sensing and continuous updating capability in the network. The essence of a digital twin is a high-precision, dynamic, and interactive digital representation of the physical world. Such a representation cannot rely only on offline modeling or static databases; it must be continuously refreshed and calibrated by persistent, reliable, and wide-area sensing mechanisms \cite{IMT2030Framework,Khan2022digital}. ISAC can utilize existing communication network infrastructure to collect information about target position, movement state, environmental structure, and scene changes while simultaneously providing connectivity, thereby supplying real-time data input for digital twins. For this reason, ISAC may be regarded as an important sensing foundation for coupling the digital world and the physical world in future networks.

Moreover, emerging scenarios such as the low-altitude economy and integrated space--air--ground--sea networking also impose stronger requirements on ``simultaneous connectivity, sensing, and decision-making'' \cite{Saad2020Vision6G,IMT2030Framework}. These scenarios are typically characterized by high dynamics, wide-area distribution, heterogeneous nodes, complex propagation environments, and rapidly changing link states, which are difficult to support adequately with conventional communication mechanisms alone. For example, in low-altitude aerial networks, the system must not only transmit control and service data, but also maintain real-time awareness of aircraft position, velocity, trajectory evolution, and surrounding airspace conditions. In integrated space--air--ground--sea networks, coordination among satellites, UAVs, terrestrial base stations, and maritime platforms also requires joint sensing and dynamic scheduling of node states, environmental conditions, and target activities. In such scenarios, ISAC enables the network to possess environmental understanding and risk warning capabilities, thereby enhancing autonomous networking, dynamic coverage, and coordinated control under complex conditions.

Overall, the 6G vision does not simply demand incremental improvement in one isolated performance metric. Rather, it requires the network to possess a comprehensive native capability oriented toward the physical world, namely the ability to sense the environment in real time, understand scenarios, predict changes, and optimize autonomously \cite{Liu2022ISAC,Gonzalez2024integrated,Saad2020Vision6G}. These requirements all point to the same development trend: future wireless networks must deeply embed sensing into their architecture, making it a foundational capability as important as connectivity itself. In this sense, ISAC is not only a key technological direction in 6G, but also one of the important foundations supporting the formation of native 6G capabilities and the transformation of network form.

\subsection{Summary of Technical Boundaries and Core Characteristics at Each Stage}
\label{subsec:ch1-evolution-summary}

After separately discussing the three-stage evolution logic of ISAC and its major driving forces, this section provides a systematic summary of the spectrum coexistence stage, the joint design stage, and the native integration stage from the three dimensions of technical objective, design level, and evaluation metric. The purpose is to help readers grasp the relationships, differences, and progressive nature of the three stages from a holistic viewpoint. It should be emphasized that this stage-based summary is not intended to artificially separate different research directions, but rather to distill the main technical thread through which the depth of communication--sensing fusion continuously increases.

\subsubsection{Differences in Technical Objectives}
\label{subsubsec:ch1-target-differences}

From the viewpoint of technical objectives, the three stages exhibit a progressive logic from ``reducing conflict,'' to ``sharing resources,'' and then to ``forming system-level closed-loop capability.''
\begin{itemize}
    \item \textbf{Stage 1: reducing conflict and maintaining compatible operation.}  
    In the spectrum coexistence stage, communication and radar/sensing systems are still usually two independent systems. The core objective at this stage is to minimize mutual interference when the two systems share the same spectrum or operate in adjacent bands, so that the performance degradation of both sides remains within acceptable limits. In other words, the goal is not to obtain functional synergy, but to ensure compatibility and coexistence. The design philosophy of this stage can thus be summarized as ``peaceful coexistence.''
    
    \item \textbf{Stage 2: sharing resources while balancing dual-function performance.}  
    In the joint design stage, the research focus shifts to ``how to let the same platform, signal, or set of resources simultaneously serve communication and sensing'' \cite{Zhang2021ISAC6G,Hassanien2015dual}. Therefore, the core objective of this stage is no longer merely conflict reduction, but rather, through dual-functional waveform design, joint beamforming, and resource sharing mechanisms, to satisfy both communication and sensing performance requirements under given constraints and to seek the optimal tradeoff between them. The design philosophy of this stage is closer to ``cooperative win--win.''
    
    \item \textbf{Stage 3: building network-level native capability and service loops.}  
    In the native integration stage, sensing becomes part of the native capability framework of 6G and future wireless networks, jointly supporting end-to-end service processes together with communication, computing, control, and intelligence \cite{Saad2020Vision6G,Zhang2020perceptive,Cui2024integrated}. Therefore, the objective is no longer limited to link-level dual-function realization, but rises further to the construction of system-level closed-loop capability for application tasks. Through the joint support of ``sensing--communication--computation--decision--control,'' this stage aims to create holistic value for scenarios such as autonomous driving, industrial control, digital twins, and low-altitude coordination. Its design philosophy may be summarized as ``native and holistic integration.''
\end{itemize}

\subsubsection{Differences in Design Levels}
\label{subsubsec:ch1-design-level-differences}

The three stages also differ in the levels at which design is focused, showing a gradual expansion from lower to higher layers and from local to global scope.
\begin{itemize}
    \item \textbf{Stage 1: spectrum and interference management layer.}  
    In the first stage, design attention is mainly concentrated on spectrum usage and interference control. The key problem is how to realize compatible operation of communication and radar systems through frequency planning, guard-band design, power control, directional isolation, cognitive spectrum access, and interference suppression \cite{Chiriyath2017RadarCommCoexistence,Blunt2018RadarSpectrumSharing}. The design object here is the \emph{resource conflict} itself, rather than a unified system architecture.
    
    \item \textbf{Stage 2: waveform, beam, and resource joint optimization layer.}  
    In the second stage, the design scope expands into physical-layer and link-layer joint design problems, including dual-functional waveform construction, MIMO joint beamforming, transmit covariance matrix optimization, and the unified scheduling of time--frequency--space--power resources \cite{Liu2020JointRadarComm,Zhang2021ISAC6G,Sturm2011WaveformDesign}. The design object has evolved from ``avoiding conflict'' to ``forming a unified signal and resource optimization framework,'' and the main technical tools include dual-functional waveform design, communication--sensing tradeoff analysis, multiobjective optimization, and constrained optimization.
    
    \item \textbf{Stage 3: network, architecture, and service layer.}  
    In the third stage, the design scope rises further to network architecture, control-plane protocols, service orchestration, and application service loops \cite{Gonzalez2024integrated,Zhang2020perceptive,Cui2024integrated}. The key problem is no longer just how a single node can communicate and sense simultaneously, but how the entire network can natively support sensing functionality, including task scheduling for multi-node cooperative sensing, sensing-data feedback and fusion, sensing management modules within network functional entities, device--edge--cloud collaborative processing, and application-oriented service interface design. At this point, ISAC has evolved from a link-level technical problem into a system-level and service-level capability construction problem.
\end{itemize}

\subsubsection{Differences in Evaluation Metrics}
\label{subsubsec:ch1-metric-differences}

As the design objectives and system levels change, the core performance metrics of interest also differ significantly across the three stages.
\begin{itemize}
    \item \textbf{Stage 1: emphasis on mutual interference suppression and compatibility evaluation.}  
    In the spectrum coexistence stage, performance evaluation mainly revolves around the question: ``How much do the performances of the two systems degrade under coexistence conditions?'' Typical metrics include interference-to-noise ratio (INR), communication SINR under interference, bit error rate, throughput, and radar/sensing metrics such as detection probability, false alarm probability, estimation error, and interference protection margin \cite{Chiriyath2017RadarCommCoexistence,Blunt2018RadarSpectrumSharing}. These metrics reflect robustness and coexistence capability in mutual-interference environments.
    
    \item \textbf{Stage 2: emphasis on communication--sensing tradeoffs.}  
    In the joint design stage, the focus shifts to ``what jointly optimal tradeoff can be achieved between communication and sensing under given resources and constraints'' \cite{Liu2020JointRadarComm,Liu2022ISAC}. Accordingly, in addition to communication-side metrics such as rate, SINR, and bit error rate, sensing-side metrics such as detection probability, estimation accuracy, CRLB, range/velocity/angle resolution, and beampattern gain are also considered. Furthermore, Pareto boundaries, rate--estimation-accuracy regions (e.g., rate--CRLB regions), and multiobjective optimal frontiers are used to characterize the performance limits of dual-functional systems \cite{Zhang2021ISAC6G}. The central question at this stage is ``how to quantify synergy gains and performance tradeoffs.''
    
    \item \textbf{Stage 3: emphasis on system-level utility and service-loop effectiveness.}  
    In the native integration stage, link-level metrics alone are no longer sufficient to fully evaluate system value. Performance evaluation must further extend to the network and service levels \cite{Zhang2020perceptive,Cui2024integrated}. In addition to the basic communication and sensing link metrics, one must consider more comprehensive indicators such as network-level sensing coverage, sensing capacity, gains from multi-node cooperative sensing, end-to-end service latency, service availability, computing and backhaul overhead, resource utilization efficiency, and the quality of closed-loop task completion. At this stage, the key evaluation question changes from ``how does a given signal or link perform?'' to ``what overall value does ISAC bring as a native network capability to system operation and application tasks?''
\end{itemize}

To provide a more intuitive comparison among the three stages, Table~\ref{tab:stage-comparison} further summarizes their major technical characteristics.

\renewcommand{\arraystretch}{1.2}
\setlength{\tabcolsep}{4pt}

\begin{longtable}{|>{\centering\arraybackslash}p{3.0cm}|
                  >{\RaggedRight\arraybackslash}p{3.8cm}|
                  >{\RaggedRight\arraybackslash}p{3.8cm}|
                  >{\RaggedRight\arraybackslash}p{3.8cm}|}
\caption{Systematic comparison of the technical characteristics of the three stages}
\label{tab:stage-comparison} \\
\hline
\textbf{Comparison Dimension} & 
\textbf{Stage 1: Spectrum Coexistence} & 
\textbf{Stage 2: Joint Design} & 
\textbf{Stage 3: Native Integration} \\
\hline
\endfirsthead

%\hline
%\textbf{Comparison Dimension} & 
%\textbf{Stage 1: Spectrum Coexistence} & 
%\textbf{Stage 2: Joint Design} & 
%\textbf{Stage 3: Native Integration} \\
%\hline
%\endhead
%
%\multicolumn{4}{r}{\small Continued on next page} \\
%\hline
%\endfoot

\endlastfoot

\textbf{Core objective} & 
Conflict mitigation & 
Resource sharing with dual functionality & 
System-level closed-loop capability \\
\hline

\textbf{Design philosophy} & 
Peaceful coexistence & 
Cooperative dual-function design & 
Native and holistic integration \\
\hline

\textbf{Design level} & 
Spectrum/interference layer & 
Waveform/beam/\allowbreak resource layer & 
Network/architecture/\allowbreak service layer \\
\hline

\textbf{System relation} & 
Coordination of two independent systems & 
Dual-functional design on one platform & 
Native capability of one network \\
\hline

\textbf{Core metrics} & 
Interference suppression and compatibility & 
Communication--sensing tradeoff & 
System-wide utility and service value \\
\hline

\textbf{Representative techniques} & 
Spectrum sharing, interference management & 
DFRC waveform, joint beamforming & 
Cellular sensing, cooperative networking, service loop \\
\hline

\textbf{Maturity level} & 
Relatively mature & 
Rapidly developing & 
Early exploration stage \\
\hline

\end{longtable}

Finally, it should be emphasized that these three stages should not be mechanically interpreted as a strictly chronological sequence. In reality, the problems of all three stages are still being pursued in parallel in both current research and engineering practice. On the one hand, spectrum coexistence remains a long-term issue in practical deployment. On the other hand, link-level joint design is still one of the central topics in present-day ISAC research, while network-level native integration is only beginning to unfold \cite{Liu2022ISAC,Gonzalez2024integrated,Cui2024integrated}. Therefore, the three-stage classification more accurately reflects a progression in the \emph{depth of technological fusion}, \emph{level of system integration}, and \emph{capability boundary}, rather than a strictly linear historical timeline. Many practical systems contain technical elements from more than one stage at the same time, which is also an important aspect of the complexity and openness of ISAC.

\subsection{Reflective Questions}
\label{subsec:ch1.1-reflect-questions}

The reflective questions in this section are divided into four categories: ``conceptual understanding,'' ``comparative analysis,'' ``in-depth thinking and discussion,'' and ``open exploration.'' Readers are encouraged to answer them independently after completing this chapter. Some of these questions may also be used for classroom discussion, group seminars, or course paper topics.

\paragraph{I. Conceptual Understanding}
\begin{enumerate}
    \item Explain in your own words why ISAC cannot be simply understood as the physical concatenation of ``communications + radar.'' Provide examples from the perspectives of the signal layer, resource layer, hardware layer, protocol layer, network layer, and service layer to illustrate the difference between ``deep coupling'' and ``simple concatenation.''

    \item In communication systems, the wireless channel is usually regarded as an unfavorable factor that must be ``overcome,'' whereas in radar/sensing systems, the channel (including target response) is the information-bearing object that must be ``analyzed.'' Based on this fundamental difference, discuss the unique challenges of channel modeling in ISAC systems: how can the same channel model simultaneously serve both communication and sensing functionalities?
    
    \begin{sloppypar}
    \item This chapter points out that there is a relationship of ``bidirectional empowerment'' between communication and sensing in ISAC. Explain, respectively, the meanings of ``sensing-assisted communication'' and ``communication-empowered sensing,'' and provide one specific application scenario for each. If either direction of this bidirectional empowerment is removed, how might system performance be affected?
    \end{sloppypar}

    \item Compare and distinguish the following terms: ISAC, JCAS, JRC, DFRC, and RadCom. Why does this textbook adopt ISAC as the unified terminology?
\end{enumerate}

\paragraph{II. Comparative Analysis}
\begin{enumerate}
    \setcounter{enumi}{4}
    \item Compare the core differences between communication systems and radar/sensing systems from the perspectives of transmit signal characteristics, receiver processing methods, target objects, performance metrics, and the role of the channel. Then summarize the main contradictions that these differences introduce into unified ISAC design.

    \item Compare the three stages of ISAC evolution---the spectrum coexistence stage, the joint design stage, and the 6G native integration stage---in terms of:
    \begin{enumerate}
        \item core technical objective;
        \item design level of concern;
        \item system relationship between communication and sensing;
        \item representative performance evaluation metrics.
    \end{enumerate}
    Then further discuss whether these three stages have a strictly chronological order. Why or why not?

    \item In the spectrum coexistence stage, representative implementation modes include time-division sharing, frequency-division isolation, coexistence-aware processing, and opportunistic spectrum reuse. Analyze the advantages and disadvantages of these four approaches, and explain why their room for improving spectral efficiency is limited, thereby motivating the evolution toward the second stage.
\end{enumerate}

\paragraph{III. In-Depth Thinking and Discussion}
\begin{enumerate}
    \setcounter{enumi}{7}
    \item This chapter points out that the core role of 5G and earlier mobile networks was mainly that of an ``information transmission pipe,'' whereas 6G networks are evolving into integrated service platforms combining connectivity, sensing, computing, and control. Based on a vertical application scenario with which you are familiar (such as autonomous driving, smart factories, or low-altitude UAV management), answer the following in detail:
    \begin{enumerate}
        \item What are the communication requirements and sensing requirements in this scenario?
        \item Why is the traditional ``external'' independent sensing solution insufficient to satisfy these requirements?
        \item How can ISAC realize the service loop of ``sensing $\rightarrow$ communication $\rightarrow$ computation $\rightarrow$ decision $\rightarrow$ control'' in this scenario?
    \end{enumerate}

    \item Communication systems operating in the mmWave/THz bands naturally possess large bandwidth and narrow beams. From the perspectives of range resolution
    \[
        \Delta R = \frac{c}{2B}
    \]
    and angular resolution, explain why high-frequency communication systems are naturally suitable for sensing. At the same time, discuss the unique technical challenges faced by high-frequency ISAC systems, such as path loss, beam alignment, hardware complexity, and energy consumption.

    \item Communication signals such as OFDM are usually random because they carry unknown data symbols, whereas traditional radar signals are usually deterministic and known in advance so that matched filtering can be performed at the receiver. In an ISAC system, if an OFDM communication waveform is used simultaneously for sensing and the sensing receiver cannot fully obtain the instantaneous transmitted data symbols, how should the receiver deal with the problem of an unknown or partially unknown reference signal? What implications does this issue have for ISAC system architecture design?

    \item This chapter summarizes the driving forces behind ISAC evolution into four aspects: spectral efficiency requirements, hardware platform development, network intelligence trends, and the 6G vision. Analyze:
    \begin{enumerate}
        \item whether there are intrinsic relationships among these four driving forces; if so, explain their logical relationship using either a diagram or text;
        \item which of these driving forces you consider the most fundamental, and why;
        \item what additional factors, beyond these four, may also promote the evolution of ISAC technology.
    \end{enumerate}

    \item In the joint beamforming design of the second stage of ISAC, the system usually needs to simultaneously satisfy the SINR requirements of communication users and the illumination power requirements in the directions of sensing targets. Analyze intuitively whether joint beamforming is easier or more difficult when the communication user and the sensing target are located in the same spatial direction. What about when they are approximately orthogonal in space? Provide your reasoning.
\end{enumerate}

\paragraph{IV. Open Exploration}
\begin{enumerate}
    \setcounter{enumi}{12}
    \item 3GPP and ITU-R have already included sensing capability in their 6G-related research agendas. Please consult the latest standardization documents or white papers (indicating the version or publication date in your answer), and answer the following:
    \begin{enumerate}
        \item What are the current scope and focal directions of ISAC standardization research?
        \item What are the major controversies or difficulties faced by ISAC in the standardization process?
        \item In your view, what key obstacles must still be overcome before ISAC can move from standardization to large-scale commercial deployment?
    \end{enumerate}

    \item This chapter points out that the sensing objects of ISAC are very broad, including not only traditional moving-target detection but also applications such as vital sign monitoring and gesture recognition. Select one nontraditional sensing application of interest to you and discuss:
    \begin{enumerate}
        \item what specific requirements this application imposes on sensing accuracy, latency, and coverage;
        \item what the basic principle is for realizing this sensing function using communication signals (such as Wi-Fi or 5G NR);
        \item what privacy and security concerns such a sensing application may raise.
    \end{enumerate}

    \item Suppose you are the CTO of a mobile operator and are evaluating whether to introduce ISAC capability into existing 5G base stations. Analyze the problem from the following perspectives:
    \begin{enumerate}
        \item What new business revenue sources might be created by introducing ISAC functionality?
        \item What hardware and software upgrades would be required for existing base stations?
        \item What technical and business risks would be involved in deploying ISAC functionality?
        \item In which scenarios would you recommend prioritizing ISAC deployment, and why?
    \end{enumerate}
\end{enumerate}

\noindent\textbf{Note:} The above reflective questions are designed in four levels: ``conceptual understanding--comparative analysis--in-depth discussion--open exploration.'' Readers are encouraged to answer them independently after studying this chapter. Some questions are also suitable as topics for classroom discussion, group seminars, or course assignments.
